\subsection{Closest methodological antecedent}
Self-collaboration Code Generation via ChatGPT~\cite{dong2023selfcollaboration} already ablates role instructions and maximum interactions in Sections~5.2--5.3. Removing roles or varying interaction is therefore not a novelty claim of the present study. Table~\ref{tab:closest-method} distinguishes the implemented interventions.

\begin{table}[htbp]\centering\small
\caption{Method comparison with the closest antecedent~\cite{dong2023selfcollaboration}, Tables~3--4 and Appendix~B.1. Interaction rounds and fresh model calls are different quantities.}\label{tab:closest-method}
\begin{tabular}{p{0.20\textwidth}p{0.31\textwidth}p{0.36\textwidth}}\toprule
Aspect & Self-collaboration paper & Component experiment\\\midrule
Labels & Role versus role-excised instructions & Identical operation sentences; only the stage prefixes change\\
Interaction & Feedback loop; maximum interaction varied & One or three fresh turns; fixed sequence, full exposed history\\
Design & Separate role and interaction ablations & Joint label-by-call matrix on every assigned task\\
Tasks & HumanEval, MBPP and extended tests & Reserved BigCodeBench sample; pre-generation reference and incorrect controls\\\bottomrule
\end{tabular}\end{table}

The joint design estimates the interaction of role labels and call topology under a fixed generation interface. The evaluated programs, native outcomes and token components make its quality and resource comparisons inspectable. This extends the empirical comparison without claiming a new principle of role collaboration.

Self-collaboration is the closest external method because its role ablation addresses the same attribution question. Our external comparison uses the pinned 2024 author implementation, whose generated-test execution differs from the simulated testing described in the 2023 paper (Section~\ref{sec:scc-feasibility}). The frozen comparison assigns SCC, fresh SR and fresh SN to the same 1,000 tasks with three repetitions; it contains no role-excised SCC arm. It therefore estimates differences between complete workflows, including feedback and stopping, while SR--SN and MR--MN isolate the specified labels within our component design. Two development checks establish execution feasibility, but the main comparison remains unfinished. D provides the minimal comparator. INoT* addresses a separate internal-dialogue algorithm question.

\subsection{Internal dialogue and repository methods}
